% !TeX program = xelatex
% !TeX root = main.tex
% 苏海川博士学位论文主文件
% 沿用用户提供的四川大学博士论文模板；兼容修复见 README.md。
\documentclass[academic,doctor,color]{./模板/scuthesis2020}

\usepackage{amsmath,amsfonts,amssymb,amsthm,mathrsfs}
\usepackage{ragged2e}
\usepackage{needspace}
% 目录中的章标题至少与其后两行同页。
\makeatletter
\AddToHook{cmd/l@chapter/before}{\Needspace{3\baselineskip}}
\makeatother
\setlength{\parskip}{3pt plus1pt minus2pt}
\setlength{\headheight}{15pt}
\allowdisplaybreaks
\setcounter{secnumdepth}{2}
\setcounter{tocdepth}{2}
% 目录初稿的空白正文占位，用于避免连续标题挤出页面；填入正文时移除。
\newcommand{\ThesisBlank}{\par\noindent\mbox{}\par}

% 基本信息集中填写；未经确认的个人信息保持空白。
% !TeX root = main.tex
% 题目为当前工作题目。导师、学号、学位类别和日期尚未确认，暂不填入。
\unitid{10610}
\STUnumber{}
\securityClassification{}
\securityYear{}
\CoverTitle{\shortstack[c]{地震应急资源服务网络韧性演化\\与协同决策研究\\{\zihao{-4}——以四川省泸定县为例}}}
\CoverSubTitle{}
\title{地震应急资源服务网络韧性演化与协同决策研究——以四川省泸定县为例}
\ENGtitle{Resilience Evolution and Coordinated Decision-Making in Earthquake Emergency Resource Service Networks: A Case Study of Luding County, Sichuan Province}
\school{灾后重建与管理学院}
\ENGschool{}
\author{苏海川}
\ENGauthor{Haichuan Su}
\supervisor{}
\ENGsupervisor{}
\degreeclass{}
\ENGdegreeclass{}
\major{安全科学与减灾}
\ENGmajor{}
\hasmajor{1}
\researchDirection{}
\direction{}
\ENGdirection{}
\completedate{}
\defensedate{}
\grantdate{}
\accomplishdate{}
\statementdate{}
\keywords{}
\ENGkeywords{}


\begin{document}
\maketitle
\frontmatter
\pagenumbering{Roman}
\pagestyle{fancy}
% !TeX root = ../main.tex
% 摘要在正文研究结果确定后撰写；此处仅保留原模板环境。
\begin{CHSabstract}
\begin{center}\heiti 摘\quad 要\end{center}
\end{CHSabstract}

\begin{ENGabstract}
\begin{center}\bfseries Abstract\end{center}
\end{ENGabstract}

\maketoc
% 正文出现图表后启用：\makefigtablist
% 符号与缩略语在统一模型确定后补充。

\mainmatter
\pagenumbering{arabic}
\pagestyle{fancy}
% !TeX root = ../main.tex
\chapter{绪论}
\label{chap:introduction}

\section{研究背景与意义}
\subsection{地震灾害下应急资源服务的连续性需求}
\subsection{道路损伤与资源供需失配的现实问题}
\subsection{服务韧性与协同决策的研究意义}

\ThesisBlank

\section{研究对象与问题界定}
\subsection{应急资源服务网络的研究边界}
\subsection{服务退化、韧性演化与协同决策问题}

\ThesisBlank

\section{研究目标与研究内容}
\subsection{总体研究目标}
\subsection{核心研究问题与研究内容}

\ThesisBlank

\section{研究方法与技术路线}
\subsection{理论建模与条件分析}
\subsection{机制实验与情景分析}
\subsection{总体技术路线}

\ThesisBlank

\section{论文结构安排}
\ThesisBlank

% !TeX root = ../main.tex
\chapter{文献综述与理论基础}
\label{chap:literature}

\section{灾害系统与网络相互依赖理论}
\subsection{灾害系统的扰动传递与级联影响}
\subsection{基础设施网络的功能依赖与服务关联}

\ThesisBlank

\section{应急资源配置与服务可达性研究}
\subsection{区域应急需求与资源能力匹配}
\subsection{道路可达性与共享资源配置}
\subsection{时限约束下的资源交付与服务保障}

\ThesisBlank

\section{服务韧性测度与动态恢复研究}
\subsection{韧性的功能表征与过程测度}
\subsection{恢复历程、需求时序与服务损失}
\subsection{动态瓶颈与恢复路径依赖}

\ThesisBlank

\section{道路状态估计与决策相关信息研究}
\subsection{稀疏观测下的道路状态估计}
\subsection{估计误差与决策损失}
\subsection{信息价值与主动观测配置}

\ThesisBlank

\section{道路恢复与资源协同决策研究}
\subsection{道路抢通排序与恢复调度}
\subsection{资源调拨与动态分配}
\subsection{互补关系、联合优化与滚动决策}

\ThesisBlank

\section{文献述评与研究切入点}
\subsection{现有研究的解释范围与证据边界}
\subsection{本文研究问题的理论定位}

\ThesisBlank

\section{本章小结}
\ThesisBlank

% !TeX root = ../main.tex
\chapter{地震应急资源服务网络建模与统一分析框架}
\label{chap:framework}

\section{应急资源服务网络的系统边界}
\subsection{道路、资源、需求与服务之间的关联}
\subsection{物资交付服务的基本分析单元}
\subsection{物理状态、观测信息与决策行动的区分}

\ThesisBlank

\section{研究区域与数据基础}
\subsection{泸定县及其区域互援联系}
\subsection{道路网络与灾害情景数据}
\subsection{人口需求与应急资源数据}

\ThesisBlank

\section{供需空间表征与参数标定}
\subsection{区域需求与资源能力的相对匹配}
\subsection{相对指标与服务模型参数的衔接}
\subsection{参数不确定性与情景分层}

\ThesisBlank

\section{道路与资源约束下的服务过程模型}
\subsection{道路通行时间、容量与路径约束}
\subsection{资源库存、在途流量与需求满足}
\subsection{需求时限与服务状态更新}

\ThesisBlank

\section{服务损失与动态韧性的统一测度}
\subsection{按时满足程度与服务缺口}
\subsection{累积服务损失与韧性轨迹}
\subsection{效率、公平与行动成本的评价边界}

\ThesisBlank

\section{理论分析与实验验证的统一设计}
\subsection{受控网络、泸定路网与灾损情景}
\subsection{共同基准、资源预算与信息条件}
\subsection{独立重复、配对比较与不确定性报告}
\subsection{机制归因、失效情景与外推边界}

\ThesisBlank

\section{本章小结}
\ThesisBlank

% !TeX root = ../main.tex
% 章边界：固定初始扰动和供需条件，分析损伤如何传导为服务损失。
% 下列标题为待研究内容，不代表命题已经证明或实验已经完成。
\chapter{地震扰动向应急资源服务退化的传导机制}
\label{chap:degradation}

\section{固定扰动条件下的应急资源服务模型}
\subsection{地震扰动与道路状态变化的情景表征}
\subsection{固定供需条件下的可行服务集}

\ThesisBlank

\section{道路损伤与服务可达性的约束传递}
\subsection{通行时间变化与时限可达性}
\subsection{容量受限与可交付资源量}

\ThesisBlank

\section{共享容量竞争与服务缺口形成}
\subsection{多需求点之间的道路容量竞争}
\subsection{共享库存约束与服务损失分配}

\ThesisBlank

\section{供需空间分布与损伤位置的交互影响}
\subsection{资源集聚与需求分散条件下的损伤响应}
\subsection{替代通道与关键供需联系}

\ThesisBlank

\section{服务退化的性质与临界条件分析}
\subsection{可行服务集的变化与服务损失边界}
\subsection{瓶颈约束激活与损失放大条件}

\ThesisBlank

\section{服务退化传导的机制实验与泸定情景分析}
\subsection{等强度不同位置扰动的受控比较}
\subsection{道路容量与供需分布的交互检验}
\subsection{约束放松下的损失归因与关键点识别}
\subsection{不同灾损情景下的传导差异与适用边界}

\ThesisBlank

\section{本章小结}
\ThesisBlank

% !TeX root = ../main.tex
% 章边界：给定外生道路恢复历程和参考配置规则，研究服务过程如何演化。
% 道路恢复行动的内生优化留给第七章。
\chapter{道路可达性与资源供需耦合下的服务韧性演化}
\label{chap:resilience}

\section{道路恢复与时变供需下的服务状态转移}
\subsection{外生恢复历程与通行能力变化}
\subsection{资源补给、需求到达与状态耦合}

\ThesisBlank

\section{需求时限与服务损失累积}
\subsection{资源等待、运输延迟与需求到期}
\subsection{阶段服务缺口与不可补偿损失}

\ThesisBlank

\section{服务韧性轨迹与演化阶段刻画}
\subsection{服务下降、受限运行与恢复过程}
\subsection{不同韧性轨迹的比较与分解}

\ThesisBlank

\section{服务恢复条件与恢复次序依赖分析}
\subsection{道路恢复与服务恢复的条件区分}
\subsection{需求到期与恢复次序的交互影响}

\ThesisBlank

\section{动态服务瓶颈辨识及其转移条件}
\subsection{道路、库存与时限约束的瓶颈判别}
\subsection{恢复阶段变化下的瓶颈转移}

\ThesisBlank

\section{服务韧性演化的机制实验与泸定情景分析}
\subsection{相同初末状态下不同恢复历程的配对比较}
\subsection{恢复历程与需求时序的交互检验}
\subsection{瓶颈变化与未恢复情景的反事实分析}
\subsection{不同恢复条件下的韧性差异与适用边界}

\ThesisBlank

\section{本章小结}
\ThesisBlank

% !TeX root = ../main.tex
% 正确模型下，无成本额外信息的最优决策价值与固定策略表现必须区分。
% 不把估计均方误差排序直接等同于服务决策损失排序。
\chapter{道路信息误差的决策影响与观测配置}
\label{chap:information}

\section{服务韧性分析中的信息约束与研究问题}
\subsection{道路状态与观测信息的区分}
\subsection{固定策略与条件最优决策的比较边界}

\ThesisBlank

\section{道路估计误差向服务损失的传递机制}
\subsection{关键路段时延高估与低估的不同影响}
\subsection{路径切换与容量瓶颈转移的阈值效应}

\ThesisBlank

\section{信息精度与决策价值的关系及成立条件}
\subsection{正确模型下无成本信息的价值边界}
\subsection{任务加权误差与策略损失的排序关系}
\subsection{观测成本与信息到达延迟的净收益影响}

\ThesisBlank

\section{面向服务决策的道路观测配置}
\subsection{观测预算与决策相关目标}
\subsection{一次性配置与序贯更新方法}
\subsection{代理目标的适用条件与可行回退}

\ThesisBlank

\section{误差传递与信息价值的机制实验}
\subsection{相同误差规模下的方向与位置对照}
\subsection{估计精度排序与决策损失排序的差异检验}
\subsection{观测配置与状态更新的收益来源分解}

\ThesisBlank

\section{泸定情景验证与适用边界}
\subsection{不同服务阶段和瓶颈状态下的信息价值}
\subsection{模型失配与误导观测的反例分析}
\subsection{观测成本和更新延迟下的收益边界}

\ThesisBlank

\section{本章小结}
\ThesisBlank

% !TeX root = ../main.tex
% 联动收益比较须使用相同预算、信息与需求条件。
% 两行动交互效应与联动策略相对分离策略的收益分别评价。
\chapter{道路抢通与资源调拨的联动决策机制}
\label{chap:coordination}

\section{道路恢复与资源交付的联动问题}
\subsection{抢通行动对通行能力的影响}
\subsection{资源调拨对服务恢复的作用}
\subsection{共同预算与行动时序约束}

\ThesisBlank

\section{抢通与调拨的统一动态决策模型}
\subsection{道路能力、库存与在途资源的状态转移}
\subsection{需求时限与服务损失评价}
\subsection{行动成本与联合可行域}

\ThesisBlank

\section{互补、替代与可分离条件分析}
\subsection{串联瓶颈下的抢通与调拨互补条件}
\subsection{替代通道与需求饱和下的边际收益变化}
\subsection{目标和约束可分解的充分条件}

\ThesisBlank

\section{联动策略与分离决策基准}
\subsection{抢通与调拨的联合规划及滚动更新}
\subsection{固定行动计划与分离规划基准}
\subsection{相同预算和信息条件下的可行性比较}

\ThesisBlank

\section{联动机制与收益来源的受控实验}
\subsection{抢通与调拨两因素交互检验}
\subsection{同预算下联动与分离策略的净收益比较}
\subsection{瓶颈位置与恢复时序的机制对照}

\ThesisBlank

\section{泸定情景验证与反例分析}
\subsection{不同灾损与资源条件下的策略表现}
\subsection{最优联动的零增益条件与近似策略的负收益来源}
\subsection{信息误差与抢通工期偏差的影响边界}

\ThesisBlank

\section{本章小结}
\ThesisBlank

% !TeX root = ../main.tex
% 结论与创新点须在研究完成后据实际证据撰写。
\chapter{结论与展望}
\label{chap:conclusions}

\section{主要研究结论}
\ThesisBlank

\section{主要创新点}
\ThesisBlank

\section{研究局限}
\ThesisBlank

\section{研究展望}
\ThesisBlank


\backmatter
% 文献条目经核验后收入 参考文献.bib；不沿用模板示例文献。
% \bibliographystyle{./模板/gbt7714-numerical}
% \bibliography{参考文献}
\chapter{参考文献}
\label{chap:references}
% 目前只建立目录，不填入未经核验的引用。
% !TeX root = ../main.tex
\chapter{攻读博士学位期间取得的研究成果}
% 根据作者、论文状态和学校要求核验后填写。

\makethanks
% 正文完成并确认送审信息后启用原模板声明页：
% \makestatement
\end{document}
